\begin{textsample}{POS Dim 1 – go – Score 37.00 – go/go\_ef\_14.txt}  \label{ex:f1_pos_108}
5.1.
Arte Pensar e \textbf{dizer} o \textbf{que} seja Arte é \textbf{uma} das \textbf{tarefas} mais complexas, na qual a \textbf{humanidade} se debruçou e vem se debruçando ao longo de sua existência.
O professor brasileiro e crítico de Arte, Jorge Coli ( [ 1995 ] 2006 ), inicia seu livro “ O \textbf{que} é arte ” afirmando \textbf{que} \textbf{inúmeros} tratados se comprometeram e se comprometem a responder à questão, tentando situá -la para definir o conceito.
Porém, as respostas ou as definições não são exclusivas, absolutas.
Pelo contrário, são divergentes, contraditórias e desenhadas pelo contexto cultural e pelo tempo histórico em \textbf{que} são produzidas.
Mesmo sem consensos, somos capazes de reconhecer algumas produções culturais como sendo arte e nos entregarmos à admiração e à interação \textbf{expressiva} e comunicativa \textbf{que} somente elas proporcionam.
Assim, “ tantas e tão diferentes são as concepções sobre a natureza da arte ” ( idem, p. 07 ).
Nesse contexto repleto de possibilidades, arriscamos a \textbf{dizer} \textbf{que} as artes são experiências \textbf{que} ampliam a \textbf{percepção} de nós mesmos, dos outros e da vida, ao possibilitar, por meio da expressão de sentimentos e emoções e da comunicação de pensamentos e ideias, os discursos poéticos, políticos, ideológicos, científicos, religiosos, por exemplo, velados e desvelados nas representações artísticas, com as quais interagimos e/ou produzimos na escola e fora dela.
Para Marcos Villela Pereira, filósofo e presidente da Federação dos Arte-Educadores do Brasil, na gestão de 1992 a 1993, as artes \textbf{configuram} \textbf{uma} área do conhecimento, pois \textbf{opera} com a organização imaginativa do \textbf{sujeito} a partir da experiência universal da \textbf{humanidade} e das experiências particulares de cada \textbf{um}, resguardados os princípios da unidade na diversidade, da harmonia na heterogeneidade e do equilíbrio nas diferenças, consolidando -se como fator de humanização ( ao resgatar a consciência da \textbf{dignidade} humana ), de socialização ( ao proporcionar a apropriação do processo criativo como \textbf{compromisso} histórico com a \textbf{humanidade} ) e de fortalecimento da identidade cultural ( gerado pela prática da experiência estética, integradora do pensar e do sentir ). ( 1993, apud PIMENTEL, 1999, p. 41 ) O campo das artes é complexo, repleto de possibilidades, \textbf{opera} e organiza a imaginação, a sensibilidade, a criatividade, a \textbf{cognição}.
Portanto, o aprender Arte na escola se torna fundamental e precisa ser \textbf{uma} experiência provocadora dos sentidos, alimentadora da experimentação, da atenção, da curiosidade, da crítica.
Também, em seus processos criativos e perceptivos, as artes precisam movimentar o \textbf{vivido} e suscitar sonhos, alimentar desejos, ressignificando, ao mesmo tempo, o estudar e o existir.
Por \textbf{que} as artes estão presentes nas escolas ocupando espaços e tempos no currículo da Educação Básica, em especial, do Ensino Fundamental?
Em contextos complexos e não lineares como os do território goiano, como as artes podem contribuir para a efetivação de \textbf{uma} educação integral, preocupada em minimizar as assimetrias culturais, sobretudo, educacionais e artísticas, \textbf{que} desenham o estado de Goiás?
Em cenários educativos formais, outrora orientados pelos Parâmetros Curriculares Nacionais ( PCN ), agora normatizados pela Base Nacional Comum Curricular ( BNCC ), como produzir experiências éticas, estéticas e poéticas sintonizadas com os desejos e necessidades dos estudantes?
Na Educação Básica, a legislação \textbf{que} organiza e orienta a escola \textbf{contemporânea} tem privilegiado as experiências com Artes Visuais, Dança, Música e Teatro.
Estas, portanto, são expressões \textbf{hegemônicas} \textbf{que} \textbf{configuram} o componente Arte, mas \textbf{que} não inviabilizam a presença e a aprendizagem de outras expressões artísticas, como as artes circenses, o audiovisual, a moda.
Em sintonia, a Base Nacional Comum Curricular ( 2017 ) normatiza e define as Artes Visuais como sendo processos e produtos artísticos e culturais, nos diversos tempos históricos e contextos sociais, \textbf{que} têm a expressão visual como elemento de comunicação.
Essas manifestações resultam de explorações plurais e transformações de materiais, de recursos tecnológicos e de apropriações da cultura cotidiana.
As Artes visuais possibilitam aos alunos explorar múltiplas culturas visuais, dialogar com as diferenças e conhecer outros espaços e possibilidades inventivas e \textbf{expressivas}, de modo a ampliar os limites escolares e criar novas formas de interação artística e de produção cultural, sejam elas concretas, sejam elas simbólicas. ( BRASIL, 2017, p. 193 ) A Música se destaca como sendo a expressão artística \textbf{que} se materializa por meio dos sons, \textbf{que} ganham forma, sentido e significado no âmbito tanto da sensibilidade subjetiva quanto das interações sociais, como resultado de saberes e valores diversos estabelecidos no domínio de cada cultura.
A ampliação e a produção dos conhecimentos musicais passam pela \textbf{percepção}, experimentação, reprodução, manipulação e criação de materiais sonoros diversos, dos mais próximos aos mais distantes da cultura musical dos alunos.
Esse processo lhes possibilita vivenciar a música \textbf{inter-relacionada} à diversidade e desenvolver saberes musicais fundamentais para sua inserção e participação crítica e ativa na sociedade. ( BRASIL, 2017, p. 194 ) Em relação ao Teatro, este instaura a experiência artística multissensorial de encontro com o outro em performance.
Nessa experiência, o corpo é \textbf{lócus} de criação ficcional de tempos, espaços e \textbf{sujeitos} distintos de si próprios, por meio do verbal, não verbal e da ação física.
Os processos de criação teatral passam por situações de criação coletiva e colaborativa, por \textbf{intermédio} de jogos, improvisações, atuações e encenações, caracterizados pela interação entre atuantes e espectadores.
O fazer teatral possibilita a intensa troca de experiências entre os alunos e \textbf{aprimora} a \textbf{percepção} estética, a imaginação, a consciência corporal, a \textbf{intuição}, a memória, a reflexão e a emoção. ( BRASIL, 2017, p. 194 ) Na sequência, afirma \textbf{que} Dança se constituiu como prática artística pelo pensamento e sentimento do corpo, mediante a articulação dos processos cognitivos e das experiências sensíveis \textbf{implicados} no movimento dançado.
Os processos de investigação e produção artística da dança centram -se naquilo \textbf{que} ocorre no e pelo corpo, discutindo e significando relações entre corporeidade e produção estética.
Ao articular os aspectos sensíveis, epistemológicos e formais do movimento dançado ao seu próprio contexto, os alunos problematizam e transformam \textbf{percepções} acerca do corpo e da dança, por meio de arranjos \textbf{que} permitem novas visões de si e do mundo.
Eles têm, assim, a oportunidade de repensar dualidades e binômios ( corpo versus \textbf{mente}, \textbf{popular} versus erudito, \textbf{teoria} versus prática ), em favor de \textbf{um} conjunto híbrido e dinâmico de práticas. ( BRASIL, 2017, p. 193 ) A experiência com qualquer \textbf{uma} dessas expressões na escola, orientada pelo Documento Curricular para Goiás - Ampliado, precisa ser \textbf{atravessada} por intencionalidades pedagógicas \textbf{que} respeitem as \textbf{singularidades} tanto de cada expressão artística quanto dos estudantes.
Por isso, deve ser realizada por profissionais especializados na área.
Desse modo, os estudantes têm oportunidade de adquirir o conhecimento, o aprofundamento e a consolidação de saberes e fazeres específicos, assim como o reconhecimento identitário, ação provisória e em permanente construção, além do desenvolvimento do sentimento de pertença cultural, tão necessário e vital.
Os currículos são resultados de escolhas didáticas, políticas, ideológicas.
Ao longo da história, as \textbf{teorias} do currículo o têm entendido como \textbf{uma} construção social \textbf{que} delimita territórios, percursos e discursos, bem como produz relações de saber, de poder e de ser.
Ou seja, o conhecimento \textbf{que} constitui o currículo está “ vitalmente, envolvido naquilo \textbf{que} somos, naquilo \textbf{que} nos tornamos : na nossa identidade, na nossa subjetividade ”, como afirma Tomaz Tadeu da Silva ( 1999, p. 15 ), professor \textbf{pesquisador} em \textbf{teoria} e currículo.
Embora, no DC-GO Ampliado, o componente curricular Arte se constitua em torno das especificidades das Artes Visuais, da Dança, da Música e do Teatro, ele apresenta pontos comuns e caros à educação na \textbf{contemporaneidade}.
Desse modo, incentiva \textbf{que} professores e estudantes interajam com as práticas artísticas/culturais relacionadas ao universo feminino, homossexual, afro-brasileiro, indígena, da cultura infanto-juvenil e dos \textbf{sujeitos} com necessidades especiais do território goiano, por exemplo, com o objetivo de ampliar as aprendizagens para além do universo masculino e europeu, \textbf{que} \textbf{historicamente} dominou os currículos da Educação Básica.
Orientada por princípios de igualdade e \textbf{equidade}, assim como éticos, políticos e estéticos \textbf{que} visam à educação integral e à construção de \textbf{uma} sociedade mais \textbf{justa}, mais democrática e mais inclusiva, a BNCC se \textbf{configurou} em \textbf{um} documento normativo e definidor das aprendizagens \textbf{que} todos os estudantes \textbf{necessitam} desenvolver ao longo das etapas e modalidades da Educação Básica.
Ou seja, ela se tornou referência nacional para os desenhos dos currículos e para o desenvolvimento das propostas político-pedagógicas das redes escolares públicas, federais, estaduais, municipais e privadas.
Espera -se, então, \textbf{que} ao longo da Educação Básica os estudantes, desafiados pela pesquisa e pela exploração, expandam seus repertórios culturais locais, regionais, nacionais e internacionais e ampliem sua imaginação, conhecimento e autonomia artística, por meio do desenvolvimento de competências.
Competências são definidas pela BNCC como sendo “ a mobilização de conhecimentos ( conceitos e procedimentos ), habilidades ( práticas cognitivas e socioemocionais ), atitudes e valores para resolver demandas complexas da vida cotidiana, do pleno exercício da cidadania e do mundo do trabalho ” ( BRASIL, 2017, p. 08 ).
Entre as dez competências gerais \textbf{que} todos os estudantes deverão desenvolver, estabelecidas pela BNCC, a \textbf{que} mais relaciona -se diretamente ao componente Arte é “ valorizar e fruir as diversas manifestações artísticas e culturais, das locais às mundiais, e também participar de práticas diversificadas da produção artístico-cultural ” ( idem, p. 09 ).
É importante ressaltar a necessidade de \textbf{inter-relacionar} as competências gerais com as competências específicas da área de linguagens e com as competências específicas do componente Arte no tratamento didático.
As competências específicas do componente Arte para o Ensino Fundamental são :

% matched lemmas: aprimorar, atravessar, cognição, compromisso, configurar, contemporaneidade, contemporâneo, dignidade, dizer, equidade, expressivo, hegemônico, historicamente, humanidade, implicar, inter-relacionar, intermédio, intuição, inúmero, justo, lócus, mente, necessitar, operar, percepção, pesquisador, popular, que, singularidade, sujeito, tarefa, teoria, um, uma, viver
\end{textsample}
